% 广西自然科学基金面上项目报告正文项目级配置。
% 项目层保留广西模板自己的字体、页边距与提纲装配；正文标题复用
% bensz-nsfc 的公共 titlesec 实现，避免在各项目重复维护标题样式。

\usepackage{xcolor}
\usepackage{graphicx}
\usepackage{amsmath}
\usepackage{amssymb}
\usepackage{bm}
\usepackage{geometry}
\usepackage{fontspec}
\usepackage{xeCJK}
\usepackage{ragged2e}
\usepackage{indentfirst}
\usepackage{enumitem}
\usepackage{caption}
\usepackage{booktabs}
\usepackage{tabularx}
\usepackage{longtable}
\usepackage{array}
\usepackage{hyperref}
\usepackage{titlesec}

% Word 基线：A4，上下 1440 twips（2.54 cm），左右 1800 twips（3.175 cm）。
\geometry{a4paper,left=3.175cm,right=3.175cm,top=2.54cm,bottom=2.54cm}

\hypersetup{
  colorlinks=true,
  urlcolor=black,
  linkcolor=black,
  anchorcolor=black,
  citecolor=black,
  CJKbookmarks=true
}

\newcommand{\GXNSFSetupFonts}{%
  \IfFileExists{bensz-fonts.sty}{%
    \usepackage{bensz-fonts}%
  }{%
    \IfFileExists{styles/bensz-fonts.sty}{%
      \makeatletter\input{styles/bensz-fonts.sty}\makeatother
      \BenszFontsSetPackageDir{styles/}%
    }{%
      \PackageError{GXNSF}{Missing required package bensz-fonts}{%
        Install bensz-fonts with the repository-level scripts/install.py before compiling the standard package.\MessageBreak
        The Overleaf release already embeds the required fonts.%
      }%
    }%
  }%
  \BenszFontsSetMainTimesFallback%
  \GXNSFPreferOriginalFonts
  \xeCJKsetup{PunctStyle=plain,CJKecglue={}}%
}

% 原 DOCX 指定“方正仿宋_GBK / 方正楷体_GBK”。优先使用系统已安装的
% 对应方正字体。若系统中没有原始字体，则由 bensz-fonts 先使用包内的
% 同款方正字体，再按 DOCX 的 Microsoft YaHei 替代记录和跨平台字体回退。
\newcommand{\GXNSFSetSystemFangsong}[1]{%
  \setCJKmainfont[AutoFakeBold=5]{#1}%
  \setCJKfamilyfont{gxnsffangsong}[AutoFakeBold=5]{#1}%
}

\newcommand{\GXNSFSetSystemKai}[1]{%
  \setCJKfamilyfont{gxnsfkai}[AutoFakeBold=5]{#1}%
}

\newcommand{\GXNSFSetFallbackFangsong}{\BenszFontsGXNSFSetupFangsongFallback}

\newcommand{\GXNSFSetFallbackKai}{\BenszFontsGXNSFSetupKaiFallback}

\newcommand{\GXNSFPreferOriginalFonts}{%
  \IfFontExistsTF{方正仿宋_GBK}{%
    \GXNSFSetSystemFangsong{方正仿宋_GBK}%
  }{%
    \IfFontExistsTF{方正仿宋简体}{%
      \GXNSFSetSystemFangsong{方正仿宋简体}%
    }{%
      \IfFontExistsTF{Microsoft YaHei}{%
        \GXNSFSetSystemFangsong{Microsoft YaHei}%
      }{\GXNSFSetFallbackFangsong}%
    }%
  }%
  \IfFontExistsTF{方正楷体_GBK}{%
    \GXNSFSetSystemKai{方正楷体_GBK}%
  }{%
    \IfFontExistsTF{方正楷体简体}{%
      \GXNSFSetSystemKai{方正楷体简体}%
    }{%
      \IfFontExistsTF{Microsoft YaHei}{%
        \GXNSFSetSystemKai{Microsoft YaHei}%
      }{\GXNSFSetFallbackKai}%
    }%
  }%
}

\GXNSFSetupFonts

\newcommand{\gxnsffangsong}{\CJKfamily{gxnsffangsong}}
\newcommand{\gxnsfkai}{\CJKfamily{gxnsfkai}}
% bensz-nsfc-headings.sty 的公共接口：只映射字体和标题参数，不引入
% bensz-nsfc-common，避免覆盖本项目独立的字号、页边距和提纲宏。
\newcommand{\templatefont}{\gxnsffangsong}
\definecolor{MsBlue}{RGB}{0,113,192}
\newcommand{\NSFCSectionBefore}{2pt plus 0pt minus 0pt}
\newcommand{\NSFCSectionAfter}{1.5pt}
\newcommand{\NSFCSubsectionIndent}{0pt}
\newcommand{\NSFCSubsectionBefore}{1.5pt plus 0pt minus 0pt}
\newcommand{\NSFCSubsectionAfter}{5pt}
\newcommand{\NSFCSubsubsectionBefore}{3pt plus 0pt minus 0pt}
\newcommand{\NSFCSubsubsectionAfter}{2pt plus 0pt minus 0pt}
\newcommand{\GXNSFFontSize}{\fontsize{16pt}{28.3pt}\selectfont}
\newcommand{\GXNSFSubheadingFontSize}{\fontsize{15pt}{26pt}\selectfont}

% 复用公共 NSFC 标题层级。每个官方提纲条目建立一个隐含的 subsection
% 编号上下文，正文一级标题沿用 Young 示例的“1.1”编号，二级标题沿用
% “（1）”编号；编号不会重复写入官方提纲文字。
\IfFileExists{bensz-nsfc-headings.sty}{%
  \RequirePackage{bensz-nsfc-headings}%
}{%
  \IfFileExists{styles/bensz-nsfc-headings.sty}{%
    \makeatletter\input{styles/bensz-nsfc-headings.sty}\makeatother
  }{%
    \PackageError{GXNSF}{Missing required package bensz-nsfc-headings}{%
      Install bensz-nsfc with the repository-level scripts/install.py before compiling.\MessageBreak
      The Overleaf release embeds this small shared headings module.%
    }%
  }%
}
\renewcommand{\subsubsectionzihao}{\GXNSFFontSize}
\newcommand{\GXNSFHeading}[1]{\subsubsection{#1}}
\newcommand{\GXNSFSubheading}[1]{\subsubsubsection{#1}}
\titleformat{\subsubsubsection}
  {\gxnsffangsong\GXNSFSubheadingFontSize\bfseries}
  {\hspace{1em} （\arabic{subsubsubsection}）}
  {0.5pt}
  {}

\pagestyle{empty}
\renewcommand{\baselinestretch}{1.0}
\setlength{\parindent}{32pt}
\setlength{\parskip}{0pt}
\XeTeXlinebreaklocale "zh"
\XeTeXlinebreakskip=0pt plus 1pt
\AtBeginDocument{\GXNSFFontSize}

\captionsetup{font={small},labelsep=period,labelfont=bf}
\setlist{nosep,leftmargin=2em}

\newcommand{\GXNSFTemplateParagraph}[2][]{%
  \par
  \begingroup
    \justifying
    \setlength{\parindent}{32pt}%
    \setlength{\parskip}{0pt}%
    \GXNSFFontSize
    #1#2\par
  \endgroup
}

\newcommand{\GXNSFReportTitle}{%
  % 锚定首页 top skip，使标题可见上边界与 Word 导出 PDF 的约 2.85 cm 对齐。
  \hrule height 0pt width 0pt\relax
  \vskip -2pt\relax
  \GXNSFTemplateParagraph[\centering\setlength{\parindent}{0pt}\gxnsffangsong\bfseries]{报告正文}%
}

\newcommand{\GXNSFPart}[1]{%
  \stepcounter{section}%
  \setcounter{subsection}{0}%
  \setcounter{subsubsection}{0}%
  \GXNSFTemplateParagraph[\gxnsfkai\raggedright\setlength{\parindent}{32pt}]{#1}%
}

% 官方提纲按 Word 基线给出固定断点；略放宽右侧可用宽度，避免加粗标题
% 增加字面宽度后把断点前的最后一个汉字挤到下一行。
\newcommand{\GXNSFItem}[2]{%
  \stepcounter{subsection}%
  \setcounter{subsubsection}{0}%
  \GXNSFTemplateParagraph[\gxnsffangsong\raggedright\setlength{\parindent}{32pt}\emergencystretch=12pt\rightskip=-12pt]{{\bfseries #1}#2}%
}

\newcommand{\GXNSFItemMixed}[3]{%
  \stepcounter{subsection}%
  \setcounter{subsubsection}{0}%
  \GXNSFTemplateParagraph[\gxnsffangsong\raggedright\setlength{\parindent}{32pt}\emergencystretch=12pt\rightskip=-12pt]{{\bfseries #1}#2{\bfseries #3}}%
}

\newcommand{\GXNSFBodyText}{%
  \setlength{\parindent}{32pt}%
  \setlength{\parskip}{0pt}%
  \GXNSFFontSize
  \gxnsffangsong
}

\pdfstringdefDisableCommands{%
  \def\gxnsffangsong{}%
  \def\gxnsfkai{}%
}
